% =====================================================================
%  Chapter 3 — Materials and Methods
%  Sections 3.1, 3.2, 3.4 and 3.8 written from the exploratory analysis.
%  Style notes: key terms in \textbf so the chapter can be skimmed;
%  enumerations use itemize rather than running prose. Citations are plain
%  numeric keys, no page or section locators.
% =====================================================================

This chapter describes the material this thesis works with and the methods
applied to it. Sections~\ref{sec:dataset} to~\ref{sec:features} establish
\textbf{what the data physically is}, how each recording was verified and how a
recording is reduced to a short list of numbers. Section~\ref{sec:stats} defines
the statistical tools used to compare the two groups. The remaining sections
cover the models and the validation protocol.

% =====================================================================
\section{Dataset}
\label{sec:dataset}

\subsection{Provenance}

Every measurement analysed in this thesis originates from a single acquisition
campaign carried out at the Universidad Politécnica de Madrid and published by
López-Delgado \textit{et al.}~\cite{lopezdelgado2026}. That study validates a
network of 24-GHz FMCW radar nodes against an optical motion-capture reference
and establishes that the network measures temporal and spatial gait parameters
accurately, including for motorically impaired participants. It stops at
\textbf{measuring gait parameters} and contains \textbf{no disease classifier},
which is precisely the gap the present work addresses.

The acquisition environment is reproduced in Figure~\ref{fig:setup_published}.
Two tripods, each carrying one trunk-height node and one ankle-height node, face
one another at the ends of a walking platform, and an eight-camera Vicon system
provides the ground truth against which the radar was
validated~\cite{lopezdelgado2026}.

\begin{figure}[H]
  \centering
  \includegraphics[width=0.72\textwidth]{paper_fig4_setup.png}
  \caption[Published acquisition environment]{The acquisition environment as
  published. Above, the corridor layout with the five phases of the
  Timed-Up-and-Go test numbered in order and the chair at the right-hand end.
  Below, a photograph of the same room. Blue marks a torso node at 1~m, green a
  foot node at 0.15~m, dark grey the processing unit, and pink a Vicon camera.
  Reproduced from~\cite{lopezdelgado2026}.}
  \label{fig:setup_published}
\end{figure}

\subsection{Acquisition protocol}

Each subject performed the \textbf{Timed-Up-and-Go} (TUG) test, a standard
clinical mobility assessment: stand up from a chair, walk 3~m in a straight
line, turn around, walk back and sit down~\cite{lopezdelgado2026}. Two
variants were recorded, three times each:

\begin{itemize}
  \item Variant \textbf{test1}, the full protocol \textbf{including the chair}, so it
        contains a sit-to-stand transition and a stand-to-sit transition;
  \item Variant \textbf{test2}, the same walk and turn \textbf{without the chair}.
\end{itemize}

\noindent This single protocol difference proves to be the most consequential
property of the dataset, for reasons developed in Section~\ref{sec:res_duration}.
Rising from a chair is among the movements most affected by Parkinson's disease,
so any group difference in elapsed time is expected to concentrate in test1.

\subsection{Cohort}

The dataset comprises \textbf{58 subjects} and \textbf{348 recordings}, split
as shown in Table~\ref{tab:cohort}. The intended design is \textbf{six
recordings per subject}: two variants by three repetitions. Fifty-six of the
58 subjects match it exactly. Two PD subjects deviate: one contributed an
extra test1 repetition (seven recordings, which is why test1 holds 76 PD
recordings rather than 75), and one is missing a test2 repetition (five
recordings, which is why test2 holds 74 rather than 75). Both cases are
examined in the quality control of Section~\ref{sec:qc} and kept.

\begin{table}[H]
  \centering
  \caption[Composition of the dataset]{Composition of the dataset.}
  \label{tab:cohort}
  \begin{tabular}{lccc}
    \toprule
    \textbf{Group} & \textbf{Subjects} & \textbf{test1 recordings} & \textbf{test2 recordings} \\
    \midrule
    Healthy control     & 33 & 99 & 99 \\
    Parkinson's disease & 25 & 76 & 74 \\
    \midrule
    \textbf{Total}      & \textbf{58} & \multicolumn{2}{c}{\textbf{348 recordings}} \\
    \bottomrule
  \end{tabular}
\end{table}

A cohort of 58 subjects is the \textbf{defining constraint of this project}. It
shapes every subsequent decision: compact models rather than large ones, reuse of
networks pretrained elsewhere rather than training from scratch, and a validation
protocol that holds out one subject at a time rather than relying on a single
fixed test split.

The dataset provides the diagnostic label, the test variant and the repetition
index. \textbf{Age, sex, disease severity and medication state are not
recorded.} Age is the most important missing item: Parkinson's disease is a
disease of later life, and ageing alone shortens step length and slows movement
\cite{mirelman2019}, so a group difference admits two explanations that the
available metadata cannot separate. Every result reported in this thesis is
therefore a difference between \textbf{these two groups as recruited}, and this
limitation is revisited in the discussion.

\subsection{Contents of one recording}

Each recording is stored as a MATLAB v5 file holding four spectrogram arrays and
two axis vectors. Only two of the arrays are used in the main analysis:

\begin{itemize}
  \item Channel \texttt{ce\_foot}, the contrast-enhanced \textbf{foot-level}
        spectrogram;
  \item Channel \texttt{ce\_torso}, the contrast-enhanced \textbf{trunk-level}
        spectrogram;
  \item A \textbf{Doppler axis} of 320 values spanning $\pm 800$~Hz and a
        \textbf{time axis} whose length varies with the duration of the walk.
\end{itemize}

\noindent Each stored channel \textbf{combines the two same-height nodes} at
the opposite ends of the corridor: at every time frame the frame from the
node with the higher signal-to-noise ratio is kept, with the far node's
Doppler axis flipped first so that both describe motion in the same
direction~\cite{lopezdelgado2026}. The files also carry the four per-node
arrays and the per-frame selection indices, which the main analysis does not
use.

\noindent The arrays arrive already oriented as Doppler against time, so no
transposition is required. Because the two enhanced spectrograms and the two axes
are the only variables read in normal operation, loading one recording costs
about $0.12$~s and the full 106~GB dataset is never read in its entirety.

\textbf{Contrast enhancement} refers to a per-recording rescaling applied by the
originating team before storage. Its purpose is straightforward: the trunk
reflects far more energy than a foot, so on an unscaled image the trunk would
saturate and the feet would be nearly invisible. The rescaling brings both into
view simultaneously. The exact operation is not published, so the enhanced
arrays are treated as the given input rather than as a reproducible step.

\subsection{Quality control}
\label{sec:qc}

All 348 recordings were verified before any analysis, so that a damaged file
could not silently become a finding. A first stage checks structure without
loading the large arrays: the expected variables are present, the four
spectrogram arrays share one shape, the Doppler axis spans the expected range,
and the duration is plausible. A second stage loads the two spectrograms and
confirms that every value is finite and that the image carries genuine
contrast.

\textbf{All 348 recordings pass every check}, so nothing was excluded and the
modest results reported later \textbf{cannot be attributed to damaged inputs}.
The two irregular subjects were examined rather than assumed: the extra
recording is a genuine repeat attempt, not a duplicate, and the missing one is
a clean deletion, so both subjects were kept.

% =====================================================================
\section{Radar System and Signal Model}
\label{sec:radar}

\subsection{From a transmitted chirp to a spectrogram}

An FMCW radar transmits a \textbf{chirp}, a brief burst whose frequency sweeps
upward, and compares the returning echo against the burst currently being
transmitted. The difference between the two encodes both the range of the
reflector and its radial velocity. Repeating this produces a continuous stream
of measurements, and a short-time Fourier transform of that stream produces the
images analysed here. Figure~\ref{fig:chain} summarises the geometry and the
processing chain.

\begin{figure}[H]
  \centering
  \includegraphics[width=\textwidth]{setup_and_signal_chain.png}
  \caption[Acquisition geometry and signal chain]{Above, the acquisition
  geometry drawn from the published parameters. Below, the five steps that turn a
  transmitted chirp into the two arrays analysed in this thesis. Steps 1 to 4
  were performed by the originating team; this work begins from the output of
  step 5.}
  \label{fig:chain}
\end{figure}

Every parameter in Table~\ref{tab:radar} was fixed by the originating team
rather than chosen for this thesis, so they are inherited constraints. They are
reproduced unchanged from~\cite{lopezdelgado2026}.

\begin{table}[H]
  \centering
  \caption[Acquisition parameters and what each one controls]{Acquisition parameters and what each one controls.}
  \label{tab:radar}
  \begin{tabular}{@{}L{3.5cm}L{2.7cm}L{8.2cm}@{}}
    \toprule
    \textbf{Parameter} & \textbf{Value} & \textbf{What it controls} \\
    \midrule
    Sweep start frequency & 23~GHz &
      The chirp sweeps upward from here, placing the system in the 24~GHz
      band; it sets the \textbf{wavelength}, about 1.3~cm, the conversion
      factor between frequency shift and velocity. \\
    Chirp bandwidth & 1.4~GHz &
      \textbf{Range resolution}: how finely two reflectors at different distances
      are separated. \\
    Chirp duration & 625~$\mu$s &
      One measurement every 625~$\mu$s is \textbf{1600 per second}, which caps the
      measurable Doppler shift at half that rate and gives the $\pm 800$~Hz span. \\
    Clutter filter & 10~Hz high-pass &
      Removes reflections from \textbf{stationary objects} such as walls and
      furniture. \\
    Analysis window & 50~ms Hann &
      The \textbf{time against frequency trade-off}; the Hann taper separates
      components roughly \textbf{30~Hz} apart (20~Hz for an ideal window). \\
    Doppler bins & 320 bins,\newline $\pm 800$~Hz &
      One stored value every 5~Hz. \\
    Radar nodes & 4 &
      Two tripods, each holding one foot-level and one trunk-level node; the
      two same-height nodes are merged into the \textbf{two stored
      channels}. \\
    Corridor / standoff & 3~m / 1~m &
      A short walk, so recordings are brief: median $8.91$~s. \\
    Beamwidth & 40$^\circ$ &
      The field of view, which sets when a subject enters and leaves the
      measurement. \\
    \bottomrule
  \end{tabular}
\end{table}

\subsection{Doppler shift and velocity}

A moving reflector shifts the frequency of the wave that returns from it. For a
monostatic radar the wave traverses the path twice, so the shift carries a factor
of two:

\begin{equation}
  f_D \;=\; \frac{2 v}{\lambda}, \qquad\text{equivalently}\qquad
  v \;=\; \frac{f_D \, \lambda}{2},
  \label{eq:doppler}
\end{equation}

\noindent where $v$ is the radial velocity of the reflector and
$\lambda = c / f_0 \approx 1.3$~cm, computed at the 23~GHz sweep start
(taking the centre of the swept band instead shifts every conversion by
about 3\,\%). A trunk advancing at $1.0$~m/s produces a shift of
approximately $153$~Hz. Equation \eqref{eq:doppler} is fixed by the
hardware, so \textbf{the vertical axis of every spectrogram can be read as a
radial-velocity axis}.

\subsection{What a single spectrogram column represents}

At any instant the radar receives reflections from \textbf{every moving part
of the body at once}, and the Fourier transform separates them by velocity. A
single spectrogram column is therefore a \textbf{distribution of the
velocities present at that moment}, not the walking speed of the subject:
brightness at height $f$ means that much of the reflected energy came from
body parts moving at the corresponding velocity.

\subsection{The two analysis bands}

Two frequency bands are used throughout. Their boundaries are inherited from the
source publication, which places them from \textbf{established walking
velocities} rather than fitting them to any particular
dataset~\cite{lopezdelgado2026}:

\begin{itemize}
  \item A \textbf{torso band} below \textbf{200~Hz}. A trunk at roughly 1~m/s
        produces about 153~Hz by Equation~\eqref{eq:doppler}, and the boundary
        allows headroom for faster walkers.
  \item A \textbf{foot band} from \textbf{200 to 500~Hz}. A foot at peak swing
        reaches roughly 3.3~m/s, which is about 506~Hz.
\end{itemize}

\noindent The data supports the split: measured on a sample of this dataset,
the torso channel concentrates its energy low and the foot channel sits
higher, as expected. One caveat is kept in mind: about $21\,\%$ of the
foot-channel energy lies above the 500~Hz boundary, so the bands
\textbf{under-report the fastest foot motion}; widening them is left for
future work.

\subsection{Structure visible in a single recording}

Three properties are directly legible from a spectrogram, and
Figure~\ref{fig:spectrograms} shows all three.

\begin{enumerate}
  \item \textbf{The turn divides every recording in two.} Energy sits below zero
        for the first half of a recording and above zero for the second. A
        negative shift means motion \textbf{away} from the node and a positive
        shift means motion \textbf{towards} it, so the sign change marks the
        instant the subject turned around. This is clinically relevant because
        turning is the slowest and least stable phase of the walk, and impaired
        turning is characteristic of Parkinson's disease.
  \item \textbf{The repeating arches are individual steps.} In the foot channel
        each arch is one foot swing, accelerating from rest, peaking mid-swing
        and decelerating to \textbf{heel strike} (foot touching ground). The valleys between arches are the
        intervals during which the foot is flat on the ground at zero velocity,
        so steps can be counted directly from the image.
  \item \textbf{The two channels capture different physics.} The foot channel is
        spiky and periodic and dominated by 200 to 500~Hz swings; the torso
        channel is a single smooth band near 0 to 200~Hz.
\end{enumerate}

\begin{figure}[H]
  \centering
  \includegraphics[width=\textwidth]{guide_data_explained.png}
  \caption[Control and PD recordings in both channels]{One control subject and
  one PD subject, each seen simultaneously by both nodes. The two recordings are
  matched for duration (8.0~s against 8.1~s) so that elapsed time is not the
  visible difference. White dashed lines mark the torso band at $\pm 200$~Hz and
  orange the foot band from 200 to 500~Hz. The two left panels are the same
  recording of the same subject viewed by the two nodes, and likewise for the two
  right panels.}
  \label{fig:spectrograms}
\end{figure}

\subsection{Coupling between the two channels}

In healthy walking the two channels are \textbf{locked to each other in time},
and the clearest expression of this is the heel strike. The source publication
marks that instant independently in each channel~\cite{lopezdelgado2026}: it
appears as a \textbf{local minimum of the foot channel's energy}, because a
foot resting on the ground is momentarily stationary, and as a \textbf{local
peak of the trunk's speed} read from the torso channel, because the trunk is
at its fastest as weight transfers. In healthy walkers the two markers coincide stride after
stride. Figure~\ref{fig:coupling} draws both markers directly on the signature of a
healthy recording of this dataset: the upper contour of the foot channel,
which dips to a valley each time the foot is planted, and the trunk-speed
trace on the torso band, which peaks at those same instants, stride after
stride.

\begin{figure}[H]
  \centering
  \includegraphics[width=\textwidth]{guide_coupling.png}
  \caption[Channel coupling at the heel strike]{The coupling made visible on
  a healthy walker. In the upper panel the solid white contour follows the
  top of the foot signature and the dashed one its lower edge; the top
  contour dips to a valley each time the foot is planted (vertical dashed
  lines). In the lower panel the dashed contours outline the torso band and
  the solid white trace is the trunk's speed, drawn at its true Doppler
  value on the band; it is near a local peak at those same instants, so the
  two channels keep time with each other.}
  \label{fig:coupling}
\end{figure}

The same publication reports the observation that motivates the modelling choices
in this thesis: that agreement is \textbf{lost in participants with Parkinson's
disease}, so the trunk no longer tracks each heel strike
reliably~\cite{lopezdelgado2026}. The
information therefore resides not in either channel alone but in
\textbf{whether the foot's rest and the trunk's speed keep time with each
other}. A measurement taken from one channel in
isolation cannot express that relationship, so any method intended to capture it
must observe both channels simultaneously.

% =====================================================================
\section{Preprocessing}
\label{sec:preproc}

Part of this chapter feeds the spectrogram itself to neural networks, and
that raises two problems: a network needs inputs of one fixed size, and 348
recordings are too few to train one. A preprocessing pipeline addresses
both, and Figure~\ref{fig:pipeoverview} shows the whole path.

\begin{figure}[H]
  \centering
  \includegraphics[width=\textwidth]{pipe_overview.png}
  \caption[The deep-learning pipeline, end to end]{The deep-learning path
  from a stored recording to one score per person. The upper row runs once
  and writes a cache to disk; the lower row runs inside every fold of the
  validation loop, so that no operation performed there can see the held-out
  subject. The classical baseline of Section~\ref{sec:baselines} follows its
  own flow and never sees windows.}
  \label{fig:pipeoverview}
\end{figure}

Two phases of the figure matter for validity. The upper row runs once and
writes a cache to disk. The lower row runs inside every round of the
validation loop and applies only the operations that must not see the
held-out subject. Throughout the rest of this chapter one such round is
called a \textbf{fold}: one subject is set aside for testing, every
data-dependent operation is performed on the remaining subjects only, and
the round is repeated with each subject in turn
(Section~\ref{sec:validation}).

Table~\ref{tab:reduction} explains each stage of
Figure~\ref{fig:pipeoverview} and traces the shape of the data through the
pipeline: a recording enters as $320 \times 12{,}000$ numbers and leaves as
one number between 0 and~1.

\begin{table}[H]
  \centering
  \caption[The stages of the deep-learning pipeline]{The stages of Figure~\ref{fig:pipeoverview}: what each one does
  and the shape it leaves behind.}
  \label{tab:reduction}
  \small
  \begin{tabular}{@{}lL{7.4cm}c@{}}
    \toprule
    Stage & What it does & Shape after it \\
    \midrule
    Windowing &
      Cuts the recording into 3.0\,s pieces, one starting every 1.5\,s
      (Section~\ref{sec:windowing}). & $320 \times 4{,}800$ \\
    Amplitude compression &
      Evens the brightness gap between the strong torso and the faint feet
      by storing $\log(1+x)$ instead of $x$. & unchanged \\
    Resizing &
      Shrinks each window to the input grid of the pretrained network. &
      $224 \times 224$ \\
    Channel stacking &
      Joins the foot and torso windows into one two-channel array. &
      $2 \times 224 \times 224$ \\
    Standardisation &
      Centres the values, using statistics from the training subjects of the
      fold only. & unchanged \\
    Augmentation &
      Training windows only: a possible flip along the Doppler axis plus a
      little noise; test windows pass unchanged. & unchanged \\
    Network &
      Scores one window: a control score and a PD score. & $2$ \\
    Window probability &
      Converts the two scores into one probability of PD. & $1$ \\
    Subject score &
      Averages one person's window probabilities: the reported score. &
      $1$ \\
    \bottomrule
  \end{tabular}
\end{table}

Several design choices in this section admit a reasonable alternative. Each
such alternative is evaluated by an \textbf{ablation}: the removal or
replacement of one single component of the pipeline, with everything else held
fixed, in order to measure how that component affects the overall result. The
choices are thereby measured rather than asserted; the ablation results are
reported in Chapter~\ref{ch:results}. One variant is computed from the raw STFT
magnitudes rather than the enhanced arrays of Section~\ref{sec:dataset}, so
the representation itself is also tested.

\subsection{Windowing}
\label{sec:windowing}

Each recording is cut into windows of \textbf{3.0\,s}, a new window starting
every \textbf{1.5\,s}, so that consecutive windows share half their content.
The length is set by the gait itself. One stride takes approximately 1.0 to
1.2\,s, and stride-to-stride regularity, one of the clinically motivated
quantities of Section~\ref{sec:features}, requires at least two complete
strides in view; a 3.0\,s window always contains two to three. The length was
fixed in advance from this reasoning rather than tuned; treating it as a
hyperparameter and selecting it with the AUC as figure of merit is left for
future work (Section~\ref{sec:futurework}). The overlap is
deliberate, for two reasons:

\begin{itemize}
  \item Advancing by a full window length would halve the number of training
        examples drawn from an already small dataset, and
  \item A stride falling across a cut would be split in two and seen whole by
        no window, whereas with a half-window overlap every moment of the walk
        appears complete in at least one window.
\end{itemize}

Figure~\ref{fig:windowing} shows the cuts drawn onto a real recording.

\begin{figure}[H]
  \centering
  \includegraphics[width=\textwidth]{pipe_windowing.png}
  \caption[Windowing of one recording]{Windowing of one real recording. Above, the cut points drawn onto
  the spectrogram; below, the resulting windows as bars, making the
  half-window overlap visible: each window starts 1.5\,s after the previous
  one and lasts 3.0\,s.}
  \label{fig:windowing}
\end{figure}

Recordings differ in duration, so they yield between 2 and 11 windows each
(4.8 on average), and the dataset of 348 recordings becomes
\textbf{1{,}673 windows}. A finer setting of 2.0\,s windows every 0.5\,s was
considered and rejected: neighbouring windows would then share 75\,\% of their
content, so the additional windows are near-copies rather than new
information.

\subsection{Fold-local standardisation}

Networks train poorly on inputs that are not centred near zero, so the
cached values are additionally \textbf{standardised} to mean~0 and standard
deviation~1. The statistics are computed from \textbf{the training
subjects of each fold only} and applied at loading time; the cache itself
holds unstandardised values. Computing the statistics over the whole dataset
would let the held-out subject influence the scale against which its own input
is later measured, a small leak, but the argument of this thesis requires
there to be none.

A common alternative standardises each window against its own mean and
spread. It is not used here, because it deletes the window's absolute
brightness, and that single deletion would remove two different things at
once:

\begin{itemize}
  \item A possible \textbf{signal}: brightness reflects how strongly the
        moving body scatters, which bradykinesia plausibly changes;
  \item A known \textbf{nuisance}: overall brightness also varies between
        recordings for reasons unrelated to gait, such as the per-recording
        contrast enhancement described in Section~\ref{sec:dataset}.
\end{itemize}

Because it cannot be known in advance which of the two dominates, the default
pipeline keeps absolute brightness and the per-window variant is evaluated as
an ablation, with the comparison reported in Chapter~\ref{ch:results}.

\subsection{Resampling and channel stacking}

Each window is resized independently from $320 \times 4{,}800$ to
$224 \times 224$ by \textbf{anti-aliased bilinear interpolation}; the target
size is not tuned but fixed by the input layer of the pretrained network of
Section~\ref{sec:deep}. Both axes remain oversampled after resizing, so the
operation discards duplication rather than detail. The two resized channels,
foot-aimed and torso-aimed, are then stacked into a single
$2 \times 224 \times 224$ array, in the same way a colour photograph stacks
its three colour layers. Figure~\ref{fig:modelinput} shows one cached window
exactly as the network receives it: the foot channel is spiky and reaches
high speeds, the torso channel is a smooth band near zero, and the
disease-relevant coupling between the two
(Section~\ref{sec:radar}) exists only when both are present.

\begin{figure}[H]
  \centering
  \includegraphics[width=\textwidth]{pipe_model_input.png}
  \caption[One cached window]{One cached window after compression, resizing and stacking:
  channel 0 is the foot-aimed node and channel 1 the torso-aimed node,
  recorded during the same three seconds of the same walk.}
  \label{fig:modelinput}
\end{figure} Each window is stored as one file, and an
\textbf{index table} records one row per window: the subject, their diagnosis,
the source recording, the position within it, and the file location. Training
operates exclusively on this index, so selecting any group of subjects
reduces to filtering the table.

\subsection{Window weighting}

Fixing the window size removes recording length from the \emph{shape} of the
input but not from the \emph{number} of inputs: longer recordings yield more
windows. PD subjects contribute 30.9 windows on average against 27.3 for
controls, and window count alone separates the groups at a subject-level AUC
of 0.62. The network never sees the count itself, but without a correction
the subjects with more windows dominate the training loss, and the
network is pushed to tune itself to whatever the longer recordings have in
common, which is exactly the group difference in duration.
The correction is a single multiplication: each window's contribution to the
training loss is scaled by the inverse of its subject's window count, so that
every subject carries exactly the same total influence regardless of how long
their recordings are. Figure~\ref{fig:windowcount} shows the shares before and
after the correction.

\begin{figure}[H]
  \centering
  \includegraphics[width=0.95\textwidth]{pipe_window_count.png}
  \caption[Training influence per subject]{Share of training influence per subject, as recorded (left) and
  after inverse-count weighting (right). Each dot is one subject; without the
  correction the shares span more than a factor of three, driven by recording
  length.}
  \label{fig:windowcount}
\end{figure}

\subsection{Direction-of-travel labelling}

The final stage attaches to every window the fraction of its foot-channel
energy located at positive Doppler frequencies, measured after subtracting the
recording's own background level. Values near 0 or 1 identify windows lying
wholly inside the outward or the return pass; values near 0.5 identify windows
containing the turn, where energy is present in both directions at once. The
label makes a targeted experiment possible. The source publication computes
its gait parameters only after discarding the turning, standing and sitting
stages~\cite{lopezdelgado2026}, whereas this pipeline keeps everything, and
the exploratory analysis of Chapter~\ref{ch:results} shows that the additional time PD
subjects take is concentrated precisely in those discarded stages. Excluding
all windows within 0.05 of the midpoint retains 54\,\% of the dataset, at the
cost of four subjects losing every window, and is evaluated as an ablation
that attacks this confound at its source. Figure~\ref{fig:turn} shows the
distribution of the label and the retention cost of excluding the turn.

\begin{figure}[H]
  \centering
  \includegraphics[width=\textwidth]{pipe_turn_v2.png}
  \caption[Direction-of-travel labelling]{Direction-of-travel labelling. Left: share of each window's
  foot-channel energy at positive Doppler; windows near 0 or 1 lie inside a
  single walking pass, windows near 0.5 contain the turn. Right: fraction of
  the dataset retained as a function of the exclusion threshold.}
  \label{fig:turn}
\end{figure}

% =====================================================================
\section{Handcrafted Features}
\label{sec:features}

A micro-Doppler signature is a grid of roughly $320 \times 12{,}000$ values. Statistical
comparison between groups requires instead a \textbf{short, fixed-length
description of each recording}, so every recording is summarised by the same
procedure: a set of descriptive statistics is computed over the whole of the
foot spectrogram, the identical calculation is repeated on the torso
spectrogram, and two cross-channel ratios are appended. The result is
\textbf{one row of 32 numbers per recording}, giving a table of 348 rows.

Of those 32, \textbf{25 are carried into the analysis}: 17 assigned to the
length-independent set and 8 to the length-dependent set
(Table~\ref{tab:features}). The
remaining 7 are redundant variants, namely the four per-channel band means,
the two alternative normalisations of the foot-to-torso ratio, and a
cross-channel ratio of the two channel means, which duplicate information
already present and are excluded to avoid inflating the number of statistical
tests.

\subsection{Behaviour with respect to recording length}

The features are divided into two classes, and the distinction governs every
later analysis. Table~\ref{tab:features} lists them.

\begin{table}[H]
  \centering
  \caption[Handcrafted features and their length behaviour]{Handcrafted features and their behaviour with recording length.}
  \label{tab:features}
  \begin{tabular}{@{}L{3.9cm}L{7.3cm}L{2.4cm}@{}}
    \toprule
    \textbf{Feature} & \textbf{Physical interpretation} & \textbf{Grows with length} \\
    \midrule
    Mean brightness       & Overall reflected energy density & no \\
    Peak brightness       & The single strongest instant of motion & no \\
    Brightness dispersion & How much the signal varies & no \\
    \textbf{Spectral centroid}  & \textbf{Mean velocity of the body}, the balance point of the energy on the Doppler axis & no \\
    \textbf{Spectral bandwidth} & \textbf{Range of velocities present}: uniform or varied motion & no \\
    Spectral entropy      & How evenly energy is spread across velocities & no \\
    \textbf{Foot-to-torso ratio} & \textbf{Share of the motion that is fast (feet) rather than slow (trunk)} & no \\
    \textbf{Rhythmicity}  & \textbf{Regularity of stride timing}, measured relative to the mean & no \\
    Band and total energy sums & Energy accumulated over the whole recording & \textbf{yes} \\
    Absolute envelope variability & Unscaled variation of energy over time & \textbf{yes} \\
    \bottomrule
  \end{tabular}
\end{table}

\subsection{Features chosen for clinical reasons}

Three features were selected because each should respond to a specific,
clinically established motor sign, and because the \textbf{direction} of the
expected response can be stated in advance. The signs themselves are taken from
the clinical literature rather than observed here: \textbf{bradykinesia} is the
one motor sign the international diagnostic criteria require before a diagnosis
of Parkinson's disease can be made~\cite{postuma2015}, while \textbf{shorter
steps with reduced foot clearance} and \textbf{increased stride-to-stride
variability} are the gait changes most consistently reported in reviews of the
condition~\cite{mirelman2019,hausdorff2009}. Table~\ref{tab:clinical} maps each
sign onto a feature.

\begin{table}[H]
  \centering
  \caption[Clinically motivated features and their predicted direction]{Clinically motivated features and their predicted direction.}
  \label{tab:clinical}
  \begin{tabular}{@{}L{3.4cm}L{4.9cm}L{6.0cm}@{}}
    \toprule
    \textbf{Motor sign} & \textbf{Effect on the walk} & \textbf{Predicted effect on the feature} \\
    \midrule
    \textbf{Bradykinesia}\newline (slowness of movement)~\cite{postuma2015} &
      Every body part moves more slowly, so all Doppler shifts are smaller. &
      The \textbf{spectral centroid falls}, because energy sits lower on the
      velocity axis. \\
    \textbf{Reduced foot clearance}\newline (shuffling, shorter steps)~\cite{mirelman2019} &
      The foot travels less far and less quickly, so less energy reaches the
      fast band. &
      The \textbf{foot-to-torso ratio falls}, because the fast share of the
      motion shrinks. \\
    \textbf{Gait variability}\newline (irregular stride to stride)~\cite{hausdorff2009} &
      Steps become less evenly timed, so energy over time is less regular. &
      \textbf{Rhythmicity rises}, because variation relative to the mean
      increases. \\
    \bottomrule
  \end{tabular}
\end{table}

Committing to these directions in advance converts each feature into a
\textbf{testable prediction} rather than an observation to be rationalised
afterwards. If the centroid rises where the disease predicts a fall, the
difference cannot be the sign the feature was chosen to detect, and something
else is driving it.

% =====================================================================
\section{Statistical Analysis}
\label{sec:stats}

The quantities defined in this section, the probability value $p$, the
effect size $\delta$ and the corrected probability value, appear in nearly
every figure and table of Chapter~\ref{ch:results}; defining them once here
keeps those captions short. All group comparisons are performed on
\textbf{one feature at a time}. For a given feature, its value is read from
every recording, producing $n_1 = 150$ values labelled PD and $n_2 = 198$
labelled control, and the three quantities below are computed from those two
lists. All three derive from the same operation: comparing every PD value
against every control value and counting how often the PD value is larger.

One caution applies to everything in this section. The tests treat the 348
recordings as independent, but each subject contributes about six, and
Chapter~\ref{ch:results} measures directly that recordings of the same
person resemble each other far more than recordings of different people.
There are therefore fewer truly independent observations than the counts
suggest, and the probability values are optimistic. This is acceptable
because of the role they play: they describe the dataset, and no conclusion
of this thesis rests on them. Every claim that matters is established by the
subject-level validation protocol of Section~\ref{sec:validation} and by
confidence intervals that resample subjects, never recordings
(Section~\ref{sec:metrics}).

\subsection{Rank-based comparison}

The \textbf{Mann-Whitney \textit{U}} statistic~\cite{mannwhitney1947} is that
count,

\begin{equation}
  U \;=\; \#\{(x_i, y_j) : x_i > y_j\} \;+\; \tfrac{1}{2}\,\#\{(x_i, y_j) : x_i = y_j\},
  \label{eq:mwu}
\end{equation}

\noindent where $x$ ranges over the PD values and $y$ over the control values.
Under the null hypothesis that the two groups are interchangeable, $U$ is
centred on $n_1 n_2 / 2$, and the \textbf{probability value} $p$ is the
probability, under that hypothesis, of a value at least as far from the centre
as the one observed, $u$:

\begin{equation}
  p \;=\; \Pr_{H_0}\!\left( \left| U - \tfrac{n_1 n_2}{2} \right| \;\ge\;
  \left| u - \tfrac{n_1 n_2}{2} \right| \right).
  \label{eq:pval}
\end{equation}

\noindent A rank-based
test was preferred to a $t$-test because several of the features are strongly
skewed, and a mean is easily displaced by a small number of extreme recordings
whereas a rank can only ever win its own pairwise comparisons.

\subsection{Effect size}

The probability value reports whether a difference is likely to be real, not
whether it is large. \textbf{Cliff's delta}~\cite{cliff1993} reports the
magnitude on a scale independent of sample size:

\begin{equation}
  \delta \;=\; \frac{\#\{x_i > y_j\} - \#\{x_i < y_j\}}{n_1 n_2}
        \;=\; \frac{2U}{n_1 n_2} - 1 .
  \label{eq:cliff}
\end{equation}

\noindent Here $\delta = 0$ means the two groups are interchangeable and
$\delta = 1$ means every PD value exceeds every control value. The magnitude
labels used throughout follow the conventional thresholds of
$0.147$, $0.33$ and $0.474$ for small, medium and large~\cite{romano2006}.
Both quantities are reported for every comparison, because a sufficiently large
dataset renders a negligible difference statistically convincing.

\subsection{Correction for multiple comparisons}

Testing 25 features in succession would produce approximately one apparently
convincing result by chance alone. Probability values are therefore corrected
with the \textbf{Benjamini-Hochberg} procedure~\cite{benjamini1995}. With $m$
tests and $p_{(1)} \le \cdots \le p_{(m)}$ the sorted probability values, the
corrected value at rank $i$ is

\begin{equation}
  p^{\mathrm{fdr}}_{(i)} \;=\; \min_{j \ge i}\;
  \min\!\left(1,\; \frac{m}{j}\, p_{(j)}\right).
  \label{eq:bh}
\end{equation}

\noindent The procedure controls the expected proportion of false positives
among the reported discoveries. It was preferred to the stricter familywise
alternative, which is valid under any dependence but so conservative on 25
strongly inter-correlated features that genuine differences would be
discarded; the Benjamini-Hochberg guarantee holds under the positive
dependence typical of such features.

\subsection{Subject-level similarity}

To quantify how identifiable an individual is from a single walk, each recording
is represented by its \textbf{17 duration-independent features}, each
standardised to zero mean and unit variance so that no feature dominates by
virtue of its numerical scale. A recording is then a point in $\mathbb{R}^{17}$
and the dissimilarity between two recordings is the Euclidean distance

\begin{equation}
  d(a, b) \;=\; \Bigg( \sum_{k=1}^{17} \big(a_k - b_k\big)^2 \Bigg)^{1/2},
  \label{eq:dist}
\end{equation}

\noindent expressed in units of one standard deviation. Restricting the
calculation to duration-independent features is deliberate: were the summed
features included, a subject who consistently walks slowly would appear
identifiable through elapsed time rather than through the shape of their gait.
For each subject two quantities are formed, the \textbf{within-subject}
distance, averaged over all pairs of that subject's own recordings, and the
\textbf{between-subject} distance, averaged from that subject's centroid to ten
recordings of other subjects drawn at random with a fixed seed.

% =====================================================================
\section{Classical Baselines}
\label{sec:baselines}

The deep models are the goal of this thesis; the classical classifiers in this
section are the \textbf{yardstick} against which they are read. Before any
network is trained, three ordinary statistical classifiers are fitted to the
25 handcrafted features of Section~\ref{sec:features}, and whatever they score
becomes the bar a network must clearly exceed before anything can be claimed
for it.

One classical prediction is produced in five stages, summarised in
Figure~\ref{fig:classflow} and repeated once for each subject in turn:

\begin{figure}[H]
  \centering
  \includegraphics[width=\textwidth]{pipe_classical_flow.png}
  \caption[One fold of the classical baseline]{One fold of the classical baseline. The loop runs once per
  subject; the 58 resulting per-subject probabilities give one AUC per
  classifier and feature set.}
  \label{fig:classflow}
\end{figure}

\begin{enumerate}
  \item \textbf{Hold out one subject.} All of that person's recordings are set
        aside and touch nothing that follows.
  \item \textbf{Fill gaps.} A guard rather than a repair: any feature
        undefined for some recording would be replaced by the median of that
        feature across the training subjects. On this dataset the guard never
        fires, since all 348 recordings yield all 25 features; the rule is
        stated because a protocol must fix in advance what happens when a
        value is missing.
  \item \textbf{Standardise.} The features span wildly different units, so
        each is rescaled to mean~0 and spread~1 using the training subjects
        only.
  \item \textbf{Fit with class weighting.} The cohort is 57\,\% control and
        43\,\% PD, so a classifier that always answers control would already
        appear 57\,\% correct. Errors on each class are therefore scaled by
        the inverse of that class's share, giving the smaller PD group the
        same total influence on the fit as the larger control group.
  \item \textbf{Score and average.} The classifier outputs a probability of
        PD, one per recording, and the held-out subject's recordings are
        averaged into \textbf{one probability per person}. Only after all 58
        folds are those 58 probabilities compared with the true diagnoses.
\end{enumerate}

Averaging per subject yields the answer a clinician would want, one per
person, and makes the score insensitive to how many recordings a subject
happens to have, the classical counterpart of the window weighting of
Section~\ref{sec:preproc}.

Three classifiers are compared (Table~\ref{tab:classifiers}), chosen not to
find a winner but because they represent different kinds of pattern, so that
their agreement or disagreement is itself informative: if only the first
succeeds, the signal is a simple trend; if only the third, it is conditional;
if all three fail together, the reasonable reading is that no signal of any of
these shapes is present at this sample size.

\begin{table}[H]
  \centering
  \caption[The three classical classifiers and their roles]{The three classical classifiers and their roles.}
  \label{tab:classifiers}
  \small
  \begin{tabular}{@{}lp{10.2cm}@{}}
    \toprule
    Model & Why it suits this problem \\
    \midrule
    Logistic regression &
      Assigns every feature a weight, adds up the evidence and converts the
      total into a probability. The most restrictive of the three, therefore
      the hardest to fool: with 58 subjects a flexible model can fit noise,
      this one largely cannot, and its weights are directly readable. \\
    \addlinespace
    Support vector machine &
      Places a curved boundary between the groups, as far as possible from
      the nearest subjects, capturing signals that appear only in
      combinations of features. \\
    \addlinespace
    Random forest &
      Averages the votes of many decision trees grown on random subsets,
      capturing thresholds and conditional patterns, and ranks the features
      by usefulness as a by-product. \\
    \bottomrule
  \end{tabular}
\end{table}

Each classifier is additionally run on five nested \textbf{feature sets},
summarised in Table~\ref{tab:featuresets}. The sets differ in exactly one
respect, how much of the recording-length confound they can contain, so
reading the resulting grid across is a controlled experiment: a score that
survives the removal of length-carrying features reflects gait, and a score
that falls step by step with them reflects recording length.

\begin{table}[H]
  \centering
  \caption[The five feature sets of the classical baseline]{The five feature sets of the classical baseline.}
  \label{tab:featuresets}
  \small
  \begin{tabular}{@{}lcp{8.2cm}@{}}
    \toprule
    Set & Size & Question it answers \\
    \midrule
    \texttt{all\_features} & 25 &
      Performance with no restraint; the number most published work would
      report. \\
    \texttt{duration\_invariant} & 17 &
      Performance once length is removed \textbf{by definition}: totals grow
      with duration, averages and ratios do not, and only the latter are
      kept. \\
    \texttt{shape\_clean} & 10 &
      Performance once length is removed \textbf{in practice}: each remaining
      feature is checked directly against measured duration, and those that
      still track it (a longer recording holds more standing and turning,
      which drags some averages) are dropped. The rule is fixed in advance:
      any feature whose rank correlation with duration reaches 0.3, computed
      once over all recordings, is out. The check uses duration only, never
      the diagnosis, and can only remove information from the classifier. \\
    \texttt{confounded\_only} & 8 &
      The totals alone; how much the confound alone contributes. \\
    \texttt{duration\_only} & 1 &
      Recording length and nothing else; \textbf{the floor}. A result not
      clearly above it demonstrates nothing about gait. \\
    \bottomrule
  \end{tabular}
\end{table}

% =====================================================================
\section{Deep Models}
\label{sec:deep}

With the yardstick defined and the input pipeline built, what remains is the
choice of networks, and that choice is dominated by one number: 58 subjects.
The 1{,}673 cached windows overstate the usable evidence, since overlapping
windows of one walk are strongly correlated, so any architecture with more
free parameters than windows can in principle memorise the training set and
learn nothing general. The two primary architectures are therefore kept
small and are deliberately chosen to bracket the problem from opposite
directions, while multi-million-parameter variants are run only as capacity
checks; a third model adapts the approach the closest literature itself
trusts most. Figure~\ref{fig:capacity} places every option against the number
of available windows.

\begin{figure}[H]
  \centering
  \includegraphics[width=0.95\textwidth]{pipe_models.png}
  \caption[Model capacity on a logarithmic scale]{Trainable parameters of each candidate on a logarithmic scale,
  against the 1{,}673 available training windows. The multi-million-parameter
  variants can memorise the dataset outright and are run only as capacity
  checks.}
  \label{fig:capacity}
\end{figure}

\subsection{SmallCNN, trained from scratch}

The first network is purpose-built and small; Figure~\ref{fig:smallcnn}
draws it layer by layer.

\begin{figure}[H]
  \centering
  \includegraphics[width=0.9\textwidth]{pipe_smallcnn.png}
  \caption[SmallCNN, layer by layer]{SmallCNN, layer by layer, with the shape of the tensor leaving
  each stage. The picture shrinks while the number of feature maps grows,
  and global average pooling replaces the flatten that would otherwise
  dominate the parameter count.}
  \label{fig:smallcnn}
\end{figure}

Three convolutional blocks of 16, 32 and 64 filters ($3 \times 3$, each
followed by batch normalisation, rectification and, for the first two,
spatial halving) feed a global average pooling layer, dropout of 0.3, and a
linear layer to two class scores, for \textbf{23{,}682 trainable parameters}
in total. Convolution recognises a local pattern wherever it occurs, which
suits gait because a stride can begin anywhere in a window. The
consequential choice is \textbf{global average pooling}: flattening the
final feature maps into the class scores would require roughly 400{,}000
parameters in that single layer, seventeen times the rest of the network,
whereas averaging removes them and forces each filter to carry meaning
regardless of where in the window it fired.

\subsection{ResNet-18 as a linear probe}

The second network brings knowledge from outside the dataset.
\textbf{Transfer learning} from ImageNet-pretrained networks is the
best-performing route on small radar datasets~\cite{park2016, ni2020}, and it
is used here in its most conservative form. Two adaptations connect the
pretrained ResNet-18 to this data. First, its input layer expects three
colour channels; the red, green and blue filters are averaged into one, and
the average is applied to both radar channels, preserving the pretrained edge
and texture detectors rather than discarding them. Second, the entire
backbone of 11.2 million weights is \textbf{frozen} at its pretrained values,
including the internal normalisation statistics, which are pinned so that
they cannot drift during training; only the final linear layer trains, with
\textbf{1{,}026 parameters}. This
configuration, a \textbf{linear probe}, asks a precise question: whether the
two groups are already separable in a rich, fixed image representation,
without allowing the network to reshape that representation around 58 people.
Figure~\ref{fig:resnet} draws the network with both adaptations marked, and
Table~\ref{tab:freeze} gives the parameter budget of each retraining depth.

\begin{figure}[H]
  \centering
  \includegraphics[width=\textwidth]{pipe_resnet.png}
  \caption[ResNet-18 as used here]{ResNet-18 as used here. Dashed stages are frozen at their
  ImageNet-pretrained values and act as a fixed description generator; only
  the final linear layer trains. The colour-filter averaging at the entrance
  is the only other modification.}
  \label{fig:resnet}
\end{figure}

\begin{table}[H]
  \centering
  \caption[Retraining depth for the pretrained network]{Retraining depth for the pretrained network.}
  \label{tab:freeze}
  \small
  \begin{tabular}{@{}lrl@{}}
    \toprule
    Setting & Trainable & Role \\
    \midrule
    Final layer only & 1{,}026 & \textbf{Default}: the linear probe. \\
    Final block and layer & 8{,}394{,}754 & Capacity check; overfitting
      expected. \\
    Everything & 11{,}174{,}402 & Not run; shown for scale. \\
    \bottomrule
  \end{tabular}
\end{table}

Two findings of the source publication shape the experiments rather than the
architecture~\cite{lopezdelgado2026} (Table~\ref{tab:papercontrib}).

\begin{table}[H]
  \centering
  \caption[What the source publication contributes to the design]{What the source publication contributes to the design.}
  \label{tab:papercontrib}
  \small
  \begin{tabular}{@{}p{5.6cm}p{6.6cm}@{}}
    \toprule
    What the paper establishes & What this design does with it \\
    \midrule
    Foot-aimed nodes remain reliable on motorically impaired subjects;
    torso-aimed nodes lose reliability on exactly those subjects. &
      Predicts the foot channel carries more of the usable signal, turning
      the foot-only run of Chapter~\ref{ch:results} into a test of a stated hypothesis
      rather than a sweep. \\
    \addlinespace
    The trunk normally accelerates slightly at every heel strike; this
    trunk-to-step coupling weakens in Parkinson's. &
      The marker lives between the channels, so both are always presented to
      the networks together. \\
    \bottomrule
  \end{tabular}
\end{table}

\subsection{Envelope-LSTM}

The third model adapts, to this dataset's combined channels, the design the
closest related work preferred over its own higher-scoring
alternative~\cite{hayashi2021}. Instead of images it
receives three velocity curves per window, each resampled to 128 steps: the
upper envelope of the foot channel, tracking the fastest speed carrying
meaningful energy, and the energy-weighted mean speeds of the foot and torso
channels. The curves feed a bidirectional LSTM with 64 hidden units per
direction, mean-pooled over time into a dropout-regularised linear layer, for
35{,}586 trainable parameters. By construction the per-recording image
texture, and any artifact living in it, is stripped away before the model
sees anything.

\subsection{Training configuration}

All deep models share the single recipe of Table~\ref{tab:trainconfig}. One
choice is physically motivated and warrants explanation: the augmentation flips
the \textbf{Doppler axis}, converting walking away from the node into the
same walk towards it, a motion every subject genuinely performs twice per
recording~\cite{fard2026}. Flipping the time axis instead was rejected as
physically wrong, since it plays a stride backwards, an
accelerate-then-brake motion no walker produces.

\begin{table}[H]
  \centering
  \caption[Shared training configuration of all deep models]{Shared training configuration of all deep models.}
  \label{tab:trainconfig}
  \small
  \begin{tabular}{@{}lp{8.6cm}@{}}
    \toprule
    Choice & Value \\
    \midrule
    Optimiser & Adam, learning rate $10^{-3}$, weight decay $10^{-4}$ \\
    Batch size & 32 windows \\
    Budget & at most 20 epochs, early stopping with patience 5 \\
    Class weighting & inverse class share, computed on effective
      (window-weighted) sample mass \\
    Window weighting & inverse window count per subject
      (Section~\ref{sec:preproc}) \\
    Augmentation & Doppler-axis flip ($p=0.5$); Gaussian noise
      ($\sigma = 0.1$ on standardised data); random masking of one span of
      time columns and one of Doppler rows \\
    \bottomrule
  \end{tabular}
\end{table}

% =====================================================================
\section{Validation Protocol}
\label{sec:validation}

Most published work in this field splits its recordings at random into
training and test sets. That choice is easier, produces far higher numbers,
and is not used here. Three findings, taken together, imply that a random
split measures the wrong thing on this dataset:

\begin{enumerate}
  \item \textbf{A walk identifies the walker.} The exploratory analysis of
        Chapter~\ref{ch:results} finds each subject's recordings roughly 2.5 times more
        similar to each other than to anyone else's, for all 58 subjects
        without exception, and published systems identify individuals from
        radar micro-Doppler at accuracies above
        93\,\%~\cite{papanastasiou2020, ni2020}. Identity is the strongest
        signal in this modality.
  \item \textbf{A random split leaks it.} With about six recordings per
        person, a random split places some of a subject's recordings in
        training and others in test, so a model can score well by recognising
        the person and recalling the label seen in training, having learned
        nothing about the disease.
  \item \textbf{The failure is documented, not hypothetical.} Of the sixteen
        works reviewed in Chapter~\ref{ch:soa}, only two report subject-independent
        results~\cite{park2016, fard2026}. In the closest study, a
        spectrogram network reached 97.8\,\% accuracy separating young from
        elderly walkers, and its own authors flagged the number as
        unreliable: the two groups had been recorded at different sites, and
        the network had learned each room's pattern of background noise, the
        ambient reflections filling the spectrogram behind the walker, rather
        than the gait~\cite{hayashi2021}.
\end{enumerate}

After the protocol had been frozen and the campaign run, these findings were
also put to a direct experimental test as a post-hoc diagnostic: the networks
of Section~\ref{sec:deep} were retrained under the literature's own random
splits, reproduced literature-grade scores, and collapsed to chance on
subjects excluded from training entirely. That experiment opens the
discussion of Chapter~\ref{ch:discussion}.

The deciding argument is simpler still. A validation protocol should match
how the tool would be used, and a screening tool is applied to people it was
never trained on. \textbf{Leave-one-subject-out} cross-validation enforces
exactly that, 58 times over. Figure~\ref{fig:split} shows the three groups of
one fold with their widths to scale, and Table~\ref{tab:split} states what
each group decides. Each fold proceeds as follows:

\begin{figure}[H]
  \centering
  \includegraphics[width=\textwidth]{pipe_split.png}
  \caption[One fold of the nested protocol]{One fold of the nested protocol, widths drawn to scale. The
  held-out subject is scored exactly once and informs no decision; the fold
  is repeated 58 times, once per subject.}
  \label{fig:split}
\end{figure}

\begin{table}[H]
  \centering
  \caption[The three groups of one fold and their roles]{The three groups of one fold and their roles.}
  \label{tab:split}
  \small
  \begin{tabular}{@{}lcp{7.4cm}@{}}
    \toprule
    Group & Size & What it decides \\
    \midrule
    Train & 49 & The weights of the network, and the standardisation
      statistics. \\
    Inner validation & 8 & Which training epoch to keep; nothing else. \\
    Test & 1 & Nothing. It is scored once and reported. \\
    \bottomrule
  \end{tabular}
\end{table}

\begin{enumerate}
  \item One subject is set aside as the \textbf{test subject}, untouched;
  \item The remaining 57 are split into 49 \textbf{training} subjects and an
        \textbf{inner validation} group of 8, drawn by a deterministic,
        stratified rule seeded by the held-out subject's identifier, so the
        split is reproducible and identical for every model; the group keeps
        close to the pool's control-to-PD proportion and always contains
        both classes, and its only job is to decide which training epoch to
        keep;
  \item The standardisation statistics of Section~\ref{sec:preproc} are
        computed from the 49 training subjects;
  \item A fresh network is trained on the training subjects' windows and,
        after every epoch, scored on the inner validation group exactly as
        test subjects are scored, one averaged probability per subject;
  \item The weights of the epoch with the highest subject-level AUC on the
        inner validation group are restored (when several epochs tie, the
        earliest wins), the test
        subject's windows are scored once, their probabilities averaged into
        one stored number, and the network is discarded;
  \item The loop restarts from a blank slate with the next subject held out.
\end{enumerate}

After 58 folds there are 58 numbers, one per subject, each produced by a
model that never saw that person, and these are compared against the 58
diagnoses to give the single reported score. The classical baselines follow
the same outer loop but omit the inner validation group: they fit in a single
pass with fixed settings (the scikit-learn defaults, version-pinned and
recorded in the frozen configuration archived with the project), so there is
no epoch to select and nothing for such a group to decide, giving 57 training
subjects per fold. The inner group is
the piece most published work omits, yet without it the only way to decide
when to stop training is to watch the test subject, which is leakage.
Leave-one-subject-out is preferred over grouped alternatives for three reasons:

\begin{itemize}
  \item \textbf{Training data is the scarce resource}, and every fold uses
        the most of it that honesty allows;
  \item It yields a \textbf{personal score for every subject}, on which the
        uncertainty and confound analyses rest;
  \item It admits \textbf{exactly one possible set of folds}, leaving nothing
        to a fortunate draw.
\end{itemize}

The complete configuration, every choice described in this chapter, was
frozen and committed to version control \textbf{before any deep model was
trained}. An adversarial review of the training code then produced a set of
corrections, committed before any re-run and informed by no test result.
Chapter~\ref{ch:results} reports every configuration that informs a conclusion and states the
total number explored, with the complete run record archived alongside the
project, because with 58 subjects the optimism introduced by selective
reporting is real and cannot be engineered away.

% =====================================================================
\section{Confound-Control Protocol}
\label{sec:confounds}

Every model, classical or deep, is read against the same set of controls,
fixed before training:

\begin{itemize}
  \item \textbf{Null floors.} Two deliberately impoverished predictors are
        scored under the identical protocol: recording duration alone, and
        the window count alone. Any model not clearly above these floors has
        not demonstrated anything about gait.
  \item \textbf{Duration controls.} The one clinically real but unspecific
        difference, PD subjects taking longer, is controlled by construction
        through the fixed-length windows of Section~\ref{sec:preproc},
        exposed by the \texttt{duration\_only} floor reported beside every
        model, and tracked by correlating each model's subject scores with
        its window counts, so that elapsed time cannot masquerade as a
        spectral finding.
  \item \textbf{Uncertainty at the subject level.} Confidence intervals are
        obtained by bootstrap resampling of \emph{subjects}, never of
        recordings or windows.
  \item \textbf{Attention verification.} Gradient-based attention maps
        (Grad-CAM) are computed for both image networks to verify which
        image regions drive their output, the direct check on shortcut
        learning that the reviewed literature lacked; they are reported
        alongside the discussion.
\end{itemize}

% =====================================================================
\section{Evaluation Metrics}
\label{sec:metrics}

Every subject leaves the validation loop with a probability of being PD.
Turning probabilities into decisions requires a threshold, and any single
threshold is a contestable choice; the headline metric therefore avoids the
choice entirely. The \textbf{area under the ROC curve} (AUC) has a direct
probabilistic reading:
\begin{equation}
  \mathrm{AUC} \;=\; P\!\left(\, s_{\mathrm{PD}} > s_{\mathrm{control}} \,\right),
  \label{eq:auc}
\end{equation}
the probability that a randomly chosen PD subject receives a higher score
than a randomly chosen control subject; 0.5 corresponds to a coin flip and
1.0 to perfect ordering. Three properties make it the right headline here:

\begin{itemize}
  \item It needs \textbf{no threshold};
  \item It is \textbf{unaffected by the 57/43 class imbalance}, which lets a
        trivial classifier reach 57\,\% accuracy;
  \item It measures \textbf{ranking}, which is what a screening tool does
        when it flags who deserves closer examination.
\end{itemize}

Balanced accuracy, sensitivity and specificity are reported alongside, with
the operating threshold set at the prevalence of the positive class rather
than 0.5, since a model trained on an imbalanced dataset is not calibrated and
a half-way cut systematically under-predicts the smaller group; the AUC
itself is independent of this choice. Finally, a caution that recurs in
Chapter~\ref{ch:results}: the literature this thesis is compared against overwhelmingly
reports \emph{accuracy} on random splits, which is not the same quantity
measured under the same conditions, and a higher number obtained that way is
not automatically the better result.
